%%%PAGE 221 rot=0 legibility=good summary=解题方法·模型化提纲（轻板模型、正交分解、极限法、往返运动、整体法、图象、平抛活用公式、椭圆运动、补偿法/等效法/图像法、计数计时除以15）
%%%BLOCK id=221-01 ch=03 sec=板块模型·轻板 kind=method alt=02 cont=none y=0.04-0.28
\textbf{1、模型化}

1.轻板
%FIG p221_a 0.10 0.05 0.37 0.1325
\notefig[0.3]{figures/p221_a.png}
$F>3\mu mg$ 时 BC 与 A 相对滑动。

%FIG p221_b 0.05 0.195 0.27 0.28
\notefig[0.25]{figures/p221_b.png}
改变$\theta$：

$\mu\le\tan\theta$ 时 A 向左下，B 向右下。 % CHECK: 不等号被涂改，难辨（≤ 或 <）

$\mu>\tan\theta$ A、B 相对静止 % CHECK: 不等号被涂改，难辨（> 或 ≥）
%%%END
%%%BLOCK id=221-02 ch=02 sec=受力分析·平衡解题方法 kind=method alt=none cont=none y=0.28-0.38
\textbf{2、}
\begin{itemize}
\item 正交分解
\item 4力\quad 先合成2定力\quad 后三力汇交 \quad 等效替代
\item 定角\quad 画力\quad 辅助圆
\end{itemize}
%%%END
%%%BLOCK id=221-03 ch=90 sec=代值·极限法 kind=method alt=none cont=none y=0.38-0.41
\textbf{3、}代值、极限法
%%%END
%%%BLOCK id=221-04 ch=01 sec=匀变速直线运动·往返运动 kind=method alt=none cont=none y=0.40-0.44
\textbf{4、}$a$恒定往返运动 $\Leftrightarrow$ 类上抛运动

↓带正负\quad $x=v_0t-\frac12at^2$ % 原稿："带正负"写在公式上方，箭头↓指向 x
%%%END
%%%BLOCK id=221-05 ch=11 sec=磁场·多解问题 kind=method alt=02,04 cont=none y=0.44-0.51
\textbf{5、}$r_{\text{轨}}$与几何长度关系不明 $\Leftrightarrow$ 多解、奇、偶

$f$的双向性
%%%END
%%%BLOCK id=221-06 ch=09 sec=库仑定律·整体法 kind=method alt=02 cont=none y=0.43-0.55
\textbf{整体法}
%FIG p221_c 0.625 0.43 0.755 0.545
\notefig[0.25]{figures/p221_c.png}
宏观视角：每个球受$\frac13mg$\unclear{向心力}

（库仑力分量）
%%%END
%%%BLOCK id=221-07 ch=02 sec=弹力·图像 kind=method alt=90 cont=none y=0.52-0.63
\textbf{6、}图象分段、曲线、轴点线斜截面
%FIG p221_d 0.055 0.55 0.20 0.62
\notefig[0.25]{figures/p221_d.png}
两个劲度系数（描点时分段了）。
%%%END
%%%BLOCK id=221-08 ch=04 sec=平抛运动·活用公式 kind=method alt=none cont=none y=0.62-0.77
\textbf{7、}活用公式
%FIG p221_e 0.04 0.64 0.25 0.765
\notefig[0.25]{figures/p221_e.png}
\begin{align*}
t_1&=\sqrt{\frac{2(h_1+h)}{g}}\\
t_2&=\sqrt{\frac{2(h_1+2h)}{g}}
\end{align*}
\[
\frac{t_1}{t_2}=\frac{\sqrt{h_1+h}}{\sqrt{h_1+2h}}<\frac{\sqrt2}{2}\ \unclear{…}\quad\text{故}\ t_2<\sqrt2\,t_1
\] % CHECK: 按算式 t1/t2 应不小于 √2/2，原稿不等号为“<”；“√2/2”后有一处涂改字迹
%%%END
%%%BLOCK id=221-09 ch=05 sec=开普勒定律·椭圆轨道 kind=method alt=none cont=none y=0.78-0.83
\textbf{8、}椭圆运动用 $v_1r_1=v_2r_2$\quad 椭圆周长$\pi a$。$\leftarrow$用动能定理\quad 时间用 $\dfrac{r_1^3}{T_1^2}=\dfrac{r_2^3}{T_2^2}$（$r_1^3$ 上方标“椭”，$r_2^3$ 上方标“圆”） % CHECK: “椭圆周长πa”难辨
%%%END
%%%BLOCK id=221-10 ch=90 sec=补偿法·等效法·图像法 kind=method alt=none cont=none y=0.84-0.96
\textbf{补偿法}
%FIG p221_f 0.175 0.84 0.30 0.89
\notefig[0.25]{figures/p221_f.png}
\unclear{$F_{\text{外大小}}=F_{\text{大小}}-F_{\text{小小}}$} % CHECK: 各下标笔迹难辨

\textbf{等效法}
%FIG p221_g 0.13 0.885 0.30 0.93
\notefig[0.25]{figures/p221_g.png}

\textbf{图像法}
%%%END
%%%BLOCK id=221-11 ch=18 sec=实验·计数计时 kind=note exp=1 alt=08 cont=none y=0.96-1.00
1到16次 \unclear{算}时 15 除以15 % 原稿："15"上方有三点及横线
%%%END
%%%PAGEEND 221
